\section{Materials and Methods}
\label{sec:methods}
\textit{Structured outline; mathematical details and complete source citations
will be developed in subsequent revisions.}

\subsection{Data sources}
We used CIMA caQTL and eQTL summary statistics (OMIX009786), cell-type-specific
pseudobulk matrices, and an external PBMC multiome WNN dataset. The MR LD
reference is hg38 1000 Genomes EAS; S-LDSC uses its separate EUR reference.
\TODO{Add study citations, donor and QTL sample sizes, published covariates,
normalization, ancestry details, and data-access statements.}

\subsection{Candidate peak--gene pairs}
Precomputed cell-type-specific candidate tables supplied the pair universe.
The documented upstream search includes peaks with centers within 1 Mb of a
gene's TSS. Current eligibility requires at least three shared QTL variants
and feature availability in both marginal datasets.
\TODO{Document the definitive candidate-table generation and coordinate rules.}

\subsection{Instrument selection and LD preparation}
The model consumes pair-indexed caQTL files documented as clumped at
\(r^2=0.99\) within 20 Mb. QTL alleles must match exactly; approximately
duplicated exposure/outcome slope pairs are removed at tolerance \(10^{-3}\).
Signed PLINK LD is aligned to the retained variants, and exposure effects,
outcome effects, and LD signs are oriented together.
\TODO{Recover the original clumping command and initial instrument-selection
thresholds; give the complete retained-variant rules.}

\subsection{GLS-IVW}
IVW is fitted through the origin after whitening by the outcome-effect
covariance constructed from outcome SEs and LD. Eigenvalue truncation targets
99.9\% variance with a five-direction minimum setting, subject to available
positive directions. The implementation uses SNP-count-based degrees of
freedom and inflates SEs only for overdispersion.
\TODO{Add the GLS formulation, reference for correlated instruments, and the
precise truncation and uncertainty conventions.}

\subsection{GLS-MR Egger}
Egger uses the same LD-aware transformation and an additional intercept
\citep{bowden2015egger}. Final tables retain the slope, SE, and p-value.
\TODO{Describe the intercept and assumptions, distinguishing the general
MR-Egger reference from this project's GLS implementation.}

\subsection{Summary-data-based Mendelian randomization}
The SMR calculation uses the retained SNP with the smallest caQTL p-value,
its outcome/exposure effect ratio, and the implemented chi-square test based
on p-derived z-scores \citep{zhu2016smr}.
\TODO{Add the exact test and SE expressions. Clarify that this implementation
does not include a HEIDI analysis.}

\subsection{Cauchy combination}
CCT combines the SMR, IVW, and Egger p-values with equal weights
\citep{liu2020cauchy}. It supplies a combined significance measure rather than
a combined causal effect estimate.

\subsection{Pseudobulk and WNN correlations}
Pseudobulk Spearman correlations use matched donors and log-transformed CPM.
WNN Spearman correlations use 80 external PBMC metacells shared across the
three analyses. Neither correlation stage includes covariate residualization.
\TODO{Complete the upstream pseudobulk and metacell construction descriptions.}

\subsection{Filtering and multiple testing}
Pairs with any missing base-method p-value are removed, then pairs with IVW
residual variance greater than 10 are removed. CCT is calculated afterward.
BH adjustment is applied separately within each gene and method
\citep{benjamini1995fdr}. Discovery uses \(\q<0.1\); GRN uses \(\q<0.05\).
The two definitions of MR-only links are stated in their Results subsections.

\subsection{Regulatory references, enrichment, and ROC-AUC}
Six reference datasets are mapped to genes and genomic intervals. Gene-level
upper-tail hypergeometric tests compare significant peaks with candidate
backgrounds; the existing global summary pools eligible gene counts and reports
fold enrichment. ROC-AUC uses \(-\log_{10}(p)\), genes with both classes, and
the existing automatic ROC direction. GTEx ROC labels additionally require
overlap with validation multiome peaks.
\TODO{Add resource citations, overlap/liftOver details, reference thresholds,
and a precise definition of the comparison population.}

\subsection{Motifs, regulatory networks, and GO}
Motif analysis uses Signac, JASPAR2020, and backgrounds matched on GC and
accessibility; enrichment requires BH-adjusted \(p<0.05\) and fold \(>1\).
TF-expression checks require detection in at least 5\% of cells. GRN overlap
uses Fisher tests against all retained tested pairs. GO BP GSEA ranks MR-only
genes by minimum MR p-value over MR-only links; a separate 5,000-draw analysis
samples tested genes scored using all their links and calculates empirical
p-values from term-specific available null NES values.
\TODO{Add resource/software citations, term-size settings, gene identifiers,
and the intended interpretation of each GO null comparison.}

\subsection{Disease heritability and GWAS overlap}
S-LDSC \citep{finucane2015sldsc} evaluates hg19 peak annotations extended by
500 bp per side, each fitted with 97 baseline-LD annotations. The nine GWAS
are listed in Figure~\ref{fig:sldsc}. A separate locus screen requires a
genome-wide-significant GWAS variant inside an unextended peak.
\TODO{Add each GWAS citation and cohort description, regression resources,
uncertainty reporting, and the variant-mapping procedure.}

\subsection{Null-SNP diagnostic}
An existing diagnostic LD-prunes variants, retains caQTL/eQTL associations
with both nominal \(p>0.05\), and examines correlations of their z-scores
(Supplementary Figure~\ref{fig:supp-qc}).
\TODO{Describe its descriptive role and distinguish SNP--pair observations
from unique variants in the pooled summary.}
